%=====================================================================
% FRONT MATTER --- title, syllabus, exam facts, how to use the volume
%=====================================================================
\pagenumbering{roman}
% Front-matter figures sit outside any chapter, so give them their own
% prefix rather than letting them inherit a misleading chapter number.
\structuralfigures{F.}
\thispagestyle{empty}

\begin{center}
\vspace*{1.1cm}

{\color{secondary}\rule{0.85\textwidth}{2pt}}\\[6pt]
{\large\mdseries\textsc{\color{secondary}Cisco Certified Network Associate}}\\[16pt]

{\Huge\bfseries\color{primary}CCNA 200-301}\\[6pt]
{\Huge\bfseries\color{primary}VERSION 2.0}\\[10pt]
{\LARGE\bfseries\color{primary}Notes, Configurations and Labs}\\[16pt]

{\color{secondary}\rule{0.85\textwidth}{2pt}}\\[16pt]

{\large\mdseries\textsc{Written from zero}}\\[6pt]
{\small\color{gray!70!black}No prior networking assumed. Protocol theory and
device configuration taught together,\\ because the exam tests both and mostly
tests the second.}\\[22pt]

% --- the five domains, announced on the title page ---
\begin{tikzpicture}[
  dm/.style={rounded corners=4pt, draw=#1, line width=1.1pt, fill=#1!7,
             text=#1, align=center, font=\scriptsize\bfseries,
             text width=2.65cm, minimum height=1.75cm}]
  \node[dm=dom1] at (-6.1,0) {\faIcon{network-wired}\\[3pt]25\%\\[1pt]
     \footnotesize\mdseries Infrastructure\\\footnotesize\mdseries and Connectivity};
  \node[dm=dom2] at (-3.05,0) {\faIcon{sitemap}\\[3pt]25\%\\[1pt]
     \footnotesize\mdseries Switching and\\\footnotesize\mdseries Network Access};
  \node[dm=dom3] at (0,0) {\faIcon{route}\\[3pt]20\%\\[1pt]
     \footnotesize\mdseries IP Routing};
  \node[dm=dom4] at (3.05,0) {\faIcon{shield-alt}\\[3pt]20\%\\[1pt]
     \footnotesize\mdseries Services and\\\footnotesize\mdseries Security};
  \node[dm=dom5] at (6.1,0) {\faIcon{robot}\\[3pt]10\%\\[1pt]
     \footnotesize\mdseries AI, Operations\\\footnotesize\mdseries and Management};
\end{tikzpicture}\\[14pt]

\begin{tcolorbox}[enhanced, width=0.88\textwidth, colback=accent!5,
  colframe=accent, boxrule=0pt, leftrule=4pt, arc=2pt, top=6pt, bottom=6pt]
\small
\textbf{\color{accent!85!black}Fifty-five per cent of this exam is hands-on.}
Of the twenty-nine topics Cisco publishes, eight begin \emph{Troubleshoot} or
\emph{Diagnose} and eight begin \emph{Configure}. Only six begin \emph{Describe}.
Under the previous blueprint, v1.1, the word ``troubleshoot'' appeared in no
topic statement at all. You cannot pass this exam by reading.
\end{tcolorbox}

\vspace{0.9cm}
{\small\color{gray!60!black}Aligned to blueprint version 2.0, live from
3 February 2027 $\bullet$ \today}
\end{center}

\newpage

%=====================================================================
\chapter*{The Exam, in One Page}
\addcontentsline{toc}{chapter}{The Exam, in One Page}
\markboth{The Exam, in One Page}{}
%=====================================================================

\begin{longtable}{@{}L{4.0cm}L{11.4cm}@{}}
\toprule
\rowcolor{primary!15}
\textbf{} & \textbf{} \\
\midrule
\endhead
\textbf{Exam} & Implementing and Administering Cisco Solutions, 200-301 CCNA,
blueprint version 2.0 \\
\rowcolor{lightbg}
\textbf{Length} & 120 minutes \\
\textbf{Live from} & 3 February 2027. Version 1.1 retires on 2 February 2027;
the exam number does not change, so check which blueprint a book or course is
written against before you trust it \\
\rowcolor{lightbg}
\textbf{Domains} & Five: Network Infrastructure and Connectivity (25\%),
Switching and Network Access (25\%), IP Routing (20\%), Network Services and
Security (20\%), AI and Network Operations and Management (10\%) \\
\textbf{Topics} & Twenty-nine, numbered 1.1 to 5.6. Every chapter of this volume
is tagged with the ones it discharges, and Appendix~\ref{app:coverage} proves all
twenty-nine are covered \\
\rowcolor{lightbg}
\textbf{The shape of it} & Eight topics ask you to \emph{troubleshoot} or
\emph{diagnose}; eight ask you to \emph{configure}; six ask you to
\emph{describe}; the rest ask you to interpret, validate, manage, select or use.
Roughly a third of the exam is fixing something broken at a command line, under
time pressure \\
\textbf{Agentic AI} & Cisco's own exam description warns that you ``may be
required to evaluate output and recommendations from agentic AI and digital
network assistants'' --- and not only in domain 5. See \chref{ai} \\
\bottomrule
\end{longtable}

\begin{alertbox}[What changed from version 1.1, and why it matters to you]
If you have an older book, it is teaching a different exam. QoS, NTP, TFTP and
FTP, the OSI and TCP/IP model descriptions, TCP versus UDP comparison,
device-role and topology-architecture theory, WLC GUI configuration, REST APIs,
JSON and Terraform have all gone. OSPFv3, hypervisors and containers, SFTP and
SCP, DNS record diagnosis, storm control, RA guard, PoE configuration, DHCP
server configuration and switch virtual interfaces have arrived --- along with a
whole domain on AI in network operations. Above all, troubleshooting went from
nothing to about a third of the paper.
\end{alertbox}

%=====================================================================
\chapter*{How to Use This Volume}
\addcontentsline{toc}{chapter}{How to Use This Volume}
\markboth{How to Use This Volume}{}
%=====================================================================

This book assumes you have never configured a network device and have not
studied networking before. It teaches the theory and the configuration together,
in the order they can actually be learned, and every chapter states which exam
topics it discharges.

\section*{The furniture in every chapter}

Each chapter is built from the same parts, always in the same order, so you
always know where you are.

\rowcolors{2}{white}{lightbg}
\begin{longtable}{@{}L{4.6cm}L{10.8cm}@{}}
\toprule
\rowcolor{primary!15}
\textbf{Element} & \textbf{What it is for} \\
\midrule
\endhead
{\scriptsize\color{neutral}\textsc{Exam topics}}~\tp{2.1} &
The blueprint topics this chapter discharges, colour-coded by domain. If you are
revising against the blueprint, navigate by these \\
\faIcon{bullseye}~\textbf{Learning Outcomes} & What you will be able to
\emph{do} after the chapter, written as actions \\
\faIcon{link}~\textbf{Before You Start} & The specific earlier chapters this one
leans on \\
\faIcon{tags}~\textbf{Key Terms} & Vocabulary introduced here. All terms also
appear in Appendix~\ref{app:glossary} \\
\faIcon{book}~\textbf{Definition} & A precise statement of a term \\
\faIcon{lightbulb}~\textbf{Concept} & The intuition behind it --- why it works
this way \\
\faIcon{terminal}~\textbf{Configure} & The actual commands, in the order you
type them. Read these with a device or simulator open, not in an armchair \\
\faIcon{search}~\textbf{Verification} & The \cmd{show} command that proves the
configuration worked, and the output you should see. Learning to read this output
is learning to troubleshoot \\
\faIcon{exclamation-circle}~\textbf{Common Pitfalls} & Mistakes made every
cohort \\
\faIcon{bug}~\textbf{Break \& Fix} & A broken device to diagnose. This is the
28\% of the exam that is troubleshooting, practised in every chapter rather than
revised once at the end \\
\faIcon{flask}~\textbf{Lab} & A hands-on exercise with a topology, an addressing
table, numbered tasks and a verification step. Labs are numbered and indexed in
Appendix~\ref{app:schedule} \\
\faIcon{graduation-cap}~\textbf{On the exam} & How this material is actually
tested, and the traps in the question wording \\
\faIcon{question-circle}~\textbf{Checkpoint} & Quick self-tests. Answers in
Appendix~\ref{app:answers} \\
\faIcon{clipboard-check}~\textbf{Chapter Summary} & The chapter compressed to
what must be retained \\
\faIcon{pen-fancy}~\textbf{Review Questions} & Exam-style questions for practice \\
\bottomrule
\end{longtable}

\section*{How to study with it}

\begin{tightnum}
  \item \textbf{Type every configuration.} Not copy and paste --- type it. The
        exam asks you to produce commands from memory under time pressure, and
        muscle memory is the only thing that survives the pressure.
  \item \textbf{Never skip the verification block.} Half of troubleshooting is
        knowing what healthy output looks like, and you only learn that by
        looking at it repeatedly when nothing is wrong.
  \item \textbf{Do the Break \& Fix before reading the answer.} Sit with the
        symptom. The exam will not hand you a diagnosis.
  \item \textbf{Build the labs, do not read them.} A lab you have read is worth
        nothing on an exam that asks you to configure.
  \item \textbf{Use Appendix~\ref{app:coverage} to revise.} It maps every one of
        the twenty-nine exam topics to the chapter that covers it, so you can
        check your own coverage rather than hoping.
\end{tightnum}

\section*{What you need}

\begin{tight}
  \item \textbf{Parts I to III: Cisco Packet Tracer.} Free with a Cisco
        Networking Academy account, runs on a modest laptop, and covers
        everything in the first three parts --- VLANs, trunking, EtherChannel,
        spanning tree, static routing and OSPF --- perfectly well.
  \item \textbf{Parts IV onwards: CML or GNS3.} From Part~IV the exam names real
        platform behaviour --- topic \tp{4.3} says \emph{IOS XE} explicitly, and
        SFTP, SCP and Ansible need a real control plane. Packet Tracer
        approximates these; it does not reproduce them.
  \item \textbf{A notebook.} Subnetting is done on paper in the exam room.
\end{tight}

\begin{alertbox}[About the configurations in this book]
Every configuration here is written to be syntactically correct for Cisco IOS
and IOS~XE, and every one follows the command order a device actually requires.
They have \textbf{not} been verified on live hardware. Platform, image and
version differences are real: a command that works on one switch may be
unavailable or spelled differently on another. Before issuing a lab to students,
run it once yourself on whatever platform your department has. That advice is
worth more than a claim of validation this book cannot honestly make.
\end{alertbox}

%=====================================================================
\chapter*{Syllabus}
\addcontentsline{toc}{chapter}{Syllabus}
\markboth{Syllabus}{}
%=====================================================================

\section*{Course description}

A first course in computer networking, taught to certification standard. It
develops the ability to design, configure, verify and troubleshoot a small
enterprise network --- switching, routing, addressing, services and security ---
and prepares the student for the Cisco CCNA 200-301 v2.0 examination. No prior
networking study is assumed. The course is built around laboratory work, because
the certification it targets is more than half practical.

\section*{Course learning outcomes}

On completion, a student will be able to:

\rowcolors{2}{white}{lightbg}
\begin{longtable}{@{}L{1.3cm}L{10.6cm}L{3.4cm}@{}}
\toprule
\rowcolor{primary!15}
\textbf{CLO} & \textbf{Outcome} & \textbf{Exam domain} \\
\midrule
\endhead
CLO-1 & Explain how data moves across a switched and routed network, and use
addressing, subnetting and the transport layer correctly. (C2, C3) &
\tp{1.3}\;\tp{1.4} \\
CLO-2 & Configure switching for a small enterprise: VLANs, trunks,
EtherChannel, spanning tree and edge-host connectivity. (C3) &
\tp{2.1}\;\tp{2.2}\;\tp{2.5} \\
CLO-3 & Configure and interpret IP routing --- static routes, single-area OSPF
for IPv4 and IPv6, and first hop redundancy. (C3, C4) &
\tp{3.1}\;\tp{3.2}\;\tp{3.3}\;\tp{3.4} \\
CLO-4 & Configure network services and security: DHCP, NAT, access control
lists, device access control and layer 2 protection. (C3) &
\tp{4.1}\;\tp{4.3}\;\tp{4.6}\;\tp{4.7} \\
CLO-5 & \textbf{Diagnose and repair a broken network} from symptoms, using
\cmd{show} output, ping, traceroute and packet capture. (C4, C5) &
\tp{2.4}\;\tp{1.6}\;\tp{1.7} \\
CLO-6 & Describe modern network operations --- management approaches, SNMP,
syslog, automation with Ansible, and the role of agentic AI --- and select an
effective prompt for an operations task. (C2, C3) &
\tp{5.1}\;\tp{5.2}\;\tp{5.3}\;\tp{5.5} \\
CLO-7 & Work safely and professionally on shared infrastructure, documenting
changes and respecting the blast radius of a configuration. (A3) & --- \\
\bottomrule
\end{longtable}

CLO-5 carries the heaviest assessment weight deliberately: it is the outcome the
exam weights most heavily and the one students are least prepared for.

\section*{Assessment}

\rowcolors{2}{white}{lightbg}
\begin{longtable}{@{}L{5.6cm}C{2.2cm}L{7.4cm}@{}}
\toprule
\rowcolor{primary!15}
\textbf{Component} & \textbf{Weight} & \textbf{Notes} \\
\midrule
\endhead
Laboratory portfolio & 30\% &
One submission per lab: topology, running configuration, and the \cmd{show}
output proving it works \\
Break \& Fix assessments & 25\% &
Timed. A pre-broken topology; find and fix the fault, and write down what it was.
Mirrors the exam's heaviest verb \\
Configuration midterm & 15\% &
At a device or simulator, not on paper \\
Theory examination & 20\% &
Written, covering the \emph{describe} and \emph{interpret} topics \\
Professional practice & 10\% &
Documentation quality, change discipline, teamwork in the lab \\
\bottomrule
\end{longtable}

A department running this course inside a degree can map the laboratory
portfolio and the Break \& Fix assessments onto a practical component, and the
remaining 35\% onto the written component, without changing the weights.

\newpage
\tableofcontents
\newpage
\listoffigures
\addcontentsline{toc}{chapter}{List of Figures}
\newpage

\pagenumbering{arabic}
\setcounter{page}{1}
